\section*{Заключение}

В работе разработан механизм автоматизированной миграции бакетов между \gls{s3}-кластерами на платформе \gls{staas}, покрывающий два режима переноса: с согласованным даунтаймом и бесшовный без остановки клиентской записи.

Решены следующие задачи:
\begin{enumerate}
    \item Проанализирован ручной процесс миграции и сформулированы требования к автоматизированному механизму.
    \item Спроектирована архитектура обоих этапов с прокси \textit{s3-shifter}, собственным мигратором (Lister/Copier) и распределёнными блокировками на объектах с TTL-арендой и fencing token.
    \item Реализованы оба этапа: даунтайм-этап~--- поверх \gls{rclone}, бесшовный~--- собственным мигратором поверх PostgreSQL и Consul.
    \item Проведено тестирование на трёх уровнях: юнит-тесты, локальный плейграунд для проверки совместимости с S3-протоколом и нагрузочное тестирование \textit{s3-shifter}.
\end{enumerate}

Дальнейшее развитие~--- адаптивное управление параллелизмом копирования вместо статического пула и расширение покрытия редких S3-сценариев (\texttt{UploadPartCopy} в смешанных backend'ах).
